\documentclass[11pt]{article}
\usepackage[margin=1in]{geometry}
\usepackage{amsmath,amssymb}
\usepackage{booktabs}
\usepackage{array}
\usepackage{hyperref}
\usepackage{xcolor}
\usepackage{enumitem}
\usepackage{longtable}
\usepackage{adjustbox}

\setlist[itemize]{leftmargin=*,topsep=2pt,itemsep=2pt}
\setlist[enumerate]{leftmargin=*,topsep=2pt,itemsep=2pt}
\newcommand{\dth}{\Delta\theta}
\newcommand{\pithirty}{$\pi/30$}
\newcommand{\pitwenty}{$\pi/20$}

\title{Phase 7 notes: comparing with optimization}
\author{}
\date{Draft notes}

\begin{document}
\maketitle


\section{Phase 7: Fourier stroke comparison}

The goal of Phase 7 is to compare the discrete strokes discovered by PPO with
smooth periodic strokes represented by a low-order temporal Fourier basis.  This
phase is motivated by the Phase 6 result that factored-action PPO with scaled
radius produces robust coherent strokes at \(N=12\), including the refined
\(\Delta\theta=\pi/40\) sweep.  Rather than further increasing \(N\) or refining
\(\Delta\theta\), we now ask whether the learned RL strokes are close to smooth
flux-optimizing loops in the same shape variables.

We write the relative joint angles as
\[
    \phi_j(\tau)
    =
    a_{j,0}
    +
    \sum_{k=1}^{K}
    \left[
        a_{j,k}\cos(2\pi k\tau)
        +
        b_{j,k}\sin(2\pi k\tau)
    \right],
    \qquad
    \tau\in[0,1],
\]
for \(j=1,\ldots,N\).  The first step is diagnostic: fit this representation to
the PPO strokes and determine how many temporal modes are needed to reproduce
the observed tip paths, cumulative spatial angles, and reward traces.  The next
step is direct optimization of the Fourier coefficients for total cycle flux,
using the same geometry, hydrodynamic model, and wall/contact constraints as in
the PPO environment.

This comparison separates three questions: whether PPO has found a smooth
low-dimensional stroke family, whether that family is close to a classical
flux-optimal loop, and how the flux-optimal stroke differs from strokes selected
by efficiency or work-penalized objectives.  Since the current PPO reward does
not include an energetic cost, we expect PPO to be closer to a flux-optimizing
stroke than to a true efficiency optimum.




\end{document}
